\chapter{Mathematical Symbols}
\raggedbottom

This page lists the symbols used in the book, grouped by topic, with
how to read each one aloud and what it means in plain words. You do not
need to learn it in advance. Come back to it whenever a formula uses a
symbol you do not recognize; the chapter where a symbol is introduced
explains it properly.

A few conventions hold throughout.
\begin{itemize}
  \item Ordinary numbers are lowercase italic letters: $x$, $a$, $b$.
  \item Vectors are lowercase \emph{bold} letters: $\vect{x}$, $\vect{w}$,
    $\vect{\mu}$.
  \item Matrices are uppercase letters: $A$, $X$, $\Sigma$.
  \item Random variables are uppercase letters, and the values they take
    are the matching lowercase letters: $X$ and $x$.
  \item Parameters that a model learns or that describe a distribution
    are usually Greek letters: $\theta$, $\mu$, $\sigma$, $\lambda$.
  \item A hat marks something estimated or predicted from data:
    $\hat\theta$, $\hat{y}$.
\end{itemize}

\newcommand{\symtable}[1]{%
  \par\smallskip
  \begin{center}\small
  \begin{tabular}{@{}>{\raggedright\arraybackslash}p{0.17\textwidth}
      >{\raggedright\arraybackslash}p{0.25\textwidth}
      >{\raggedright\arraybackslash}p{0.52\textwidth}@{}}
    \toprule
    Symbol & Read as & Meaning \\
    \midrule
    #1
    \bottomrule
  \end{tabular}
  \end{center}}

\section*{Numbers, Sets, and Comparisons}

\symtable{
  $\R$ & ``the reals'' & All ordinary numbers, positive, negative, and
    zero, including fractions and decimals. \\
  $\R^n$ & ``R n'' & All lists (vectors) of $n$ real numbers. \\
  $x \in S$ & ``$x$ in $S$'' & $x$ is one of the members of the set
    $S$. \\
  $S \subset T$ & ``$S$ subset of $T$'' & Every member of $S$ is also in
    $T$. \\
  $\{1, 2, 3\}$ & ``the set 1, 2, 3'' & A set, listed by its members. \\
  $[a, b]$, $(a, b)$ & ``closed / open interval'' & All numbers from $a$ to
    $b$; square brackets include the ends, round ones leave them out. \\
  $\infty$ & ``infinity'' & Larger than every number; $(0, \infty)$ means
    all positive numbers. \\
  $|x|$ & ``absolute value of $x$'' & The size of $x$ without its sign:
    $|{-3}| = 3$. \\
  $=$, $\ne$, $\approx$ & ``equals, not equal, approximately'' & Exactly
    the same, different, close enough for the purpose. \\
  $<$, $\le$, $>$, $\ge$ & ``less than, at most, \dots'' & Comparisons;
    the line underneath allows equality. \\
  $\propto$ & ``proportional to'' & Equal up to a constant factor that we
    do not care about. \\
  $\Rightarrow$, $\Leftrightarrow$ & ``implies, if and only if'' & The left
    side leads to the right; the two sides say the same thing. \\
  $n!$ & ``$n$ factorial'' & $1 \times 2 \times \dots \times n$; for example,
    $4! = 24$. \\
  $\binom{n}{k}$ & ``$n$ choose $k$'' & The number of ways to pick $k$ items
    out of $n$. \\
  $\Delta_n$ & ``the simplex'' & All vectors of $n$ nonnegative numbers
    that add up to $1$: the valid probability vectors. \\
}

\section*{Functions}

\symtable{
  $f(x)$ & ``$f$ of $x$'' & The output of the function $f$ for the input
    $x$. \\
  $f\colon A \to B$ & ``$f$ from $A$ to $B$'' & $f$ takes inputs from the
    set $A$ and returns outputs in the set $B$. \\
  $x \mapsto x^2$ & ``$x$ maps to $x$ squared'' & A function written by
    showing what it does to an input. \\
  $\dom f$ & ``domain of $f$'' & The inputs for which $f$ is defined. \\
  $e^x$, $\exp(x)$ & ``e to the $x$'' & The exponential function, with $e
    \approx 2.718$. \\
  $\log x$ & ``log $x$'' & The natural logarithm, the inverse of $e^x$:
    $\log e^x = x$. \\
  $\max$, $\min$ & ``max, min'' & The largest or smallest value. \\
  $\operatorname*{arg\,max}$, $\operatorname*{arg\,min}$ & ``arg max, arg
    min'' & The \emph{input} at which the largest or smallest value is
    reached, not the value itself. \\
  $x^\star$ & ``$x$ star'' & The best input, usually a minimizer. \\
  $\sigma(z)$ & ``sigmoid of $z$'' & $1/(1 + e^{-z})$: squashes any number
    into $(0, 1)$. \\
  $\operatorname{ReLU}(z)$ & ``ReLU'' & $\max(0, z)$: keeps positive
    numbers, turns negative ones into $0$. \\
  $\operatorname{softmax}(\vect{z})$ & ``softmax'' & Turns a vector of
    scores into a probability vector. \\
  $\operatorname{sign}(x)$ & ``sign of $x$'' & $1$ if $x > 0$, $-1$ if $x
    < 0$, $0$ if $x = 0$. \\
}

\section*{Sums, Products, and Calculus}

\symtable{
  $\sum_{i=1}^{n} a_i$ & ``sum from $i = 1$ to $n$'' & $a_1 + a_2 + \dots +
    a_n$. \\
  $\prod_{i=1}^{n} a_i$ & ``product from $i = 1$ to $n$'' & $a_1 \times a_2
    \times \dots \times a_n$. \\
  $\lim_{h \to 0}$ & ``limit as $h$ goes to $0$'' & The value something
    approaches as $h$ gets closer and closer to $0$. \\
  $f'(x)$ & ``$f$ prime of $x$'' & The derivative: the slope of $f$ at $x$,
    how fast $f$ changes. \\
  $f''(x)$ & ``$f$ double prime'' & The second derivative: how fast the
    slope changes, which measures curvature. \\
  $\dfrac{\partial f}{\partial x_i}$ & ``partial $f$, partial $x_i$'' & The
    slope of $f$ along one input $x_i$, holding the others fixed. \\
  $\nabla f$ & ``gradient of $f$'' & The vector of all partial
    derivatives; it points in the direction of steepest increase. \\
  $\int_a^b f(x)\,dx$ & ``integral from $a$ to $b$'' & The area under the
    graph of $f$ between $a$ and $b$. \\
  $\eta$ & ``eta'' & The learning rate, the step size of gradient
    descent. \\
  $x_k$ & ``$x$ sub $k$'' & In an algorithm, the value after $k$ steps. \\
  $w \leftarrow w - \eta g$ & ``$w$ becomes \dots'' & Update: replace $w$ by
    the value on the right. \\
}

\section*{Vectors and Matrices}

\symtable{
  $\vect{x}$ & ``vector $x$'' & A list of numbers, $(x_1, \dots, x_n)$,
    treated as a column. \\
  $x_i$ & ``$x$ sub $i$'' & The $i$th entry of the vector $\vect{x}$. \\
  $\vect{x}_i$ & ``$x$ sub $i$'' (bold) & The $i$th vector in a list, such as
    the $i$th data point. \\
  $A$, $a_{ij}$ & ``$A$; $a$ sub $i j$'' & A matrix, a grid of numbers; the
    entry in row $i$, column $j$. \\
  $A^\top$, $\vect{x}^\top$ & ``$A$ transpose'' & Rows and columns swapped;
    $\vect{x}^\top$ is $\vect{x}$ written as a row. \\
  $\vect{x}^\top\vect{y}$, $\vect{x} \cdot \vect{y}$, $\langle \vect{x},
    \vect{y} \rangle$ & ``$x$ dot $y$'' & The dot product: multiply
    matching entries and add. Measures how aligned two vectors are. \\
  $\|\vect{x}\|$ & ``norm of $x$'' & The length of $\vect{x}$,
    $\sqrt{x_1^2 + \dots + x_n^2}$. \\
  $I$ & ``the identity'' & The matrix that leaves every vector unchanged:
    $1$s on the diagonal, $0$s elsewhere. \\
  $A^{-1}$ & ``$A$ inverse'' & The matrix that undoes $A$:
    $A^{-1}A = I$. \\
  $\det A$ & ``determinant of $A$'' & How much $A$ scales area or volume;
    $0$ means $A$ squashes space flat. \\
  $\operatorname{tr} A$ & ``trace of $A$'' & The sum of the diagonal
    entries. \\
  $\operatorname{diag}(\dots)$ & ``diag'' & A matrix with the given numbers
    on the diagonal and $0$s elsewhere. \\
  $A \odot B$ & ``$A$ Hadamard $B$'' & Multiply matching entries, one by
    one. \\
  $\operatorname{proj}$ & ``projection'' & The shadow of one vector on a
    line or plane. \\
  $\lambda$, $\vect{v}$ & ``lambda; $v$'' & An eigenvalue and eigenvector:
    $A\vect{v} = \lambda\vect{v}$, so $A$ only stretches $\vect{v}$. \\
  $Q\Lambda Q^\top$ & ``$Q$ Lambda $Q$ transpose'' & Eigendecomposition of
    a symmetric matrix: rotate, stretch, rotate back. \\
  $U\Sigma V^\top$, $\sigma_i$ & ``$U$ Sigma $V$ transpose'' & The singular
    value decomposition; the $\sigma_i$ are the singular values, the
    stretch factors. \\
  $X$, $\tilde{X}$ & ``$X$; $X$ tilde'' & The data matrix, one data point
    per row; the same data with the mean subtracted. \\
}

\section*{Probability}

\symtable{
  $\Omega$, $\omega$ & ``omega'' & The sample space, all possible
    outcomes; one outcome. \\
  $P(A)$ & ``probability of $A$'' & How likely the event $A$ is, from $0$
    (never) to $1$ (always). \\
  $A \cap B$ & ``$A$ and $B$'' & The event that both $A$ and $B$
    happen. \\
  $A^c$ & ``not $A$'' & The event that $A$ does not happen. \\
  $P(A \mid B)$ & ``probability of $A$ given $B$'' & How likely $A$ is once
    we know that $B$ happened. \\
  $X$, $x$ & ``big $X$, small $x$'' & A random variable, and a particular
    value it can take. \\
  $p(x)$ & ``p of $x$'' & The probability of the value $x$ (discrete), or
    the density at $x$ (continuous). \\
  $p(x; \theta)$ & ``p of $x$ with parameter $\theta$'' & A distribution
    that depends on the parameter $\theta$. \\
  $p(x \mid y)$ & ``p of $x$ given $y$'' & The distribution of $x$ once
    $y$ is known. \\
  $F_X(x)$ & ``CDF of $X$'' & $P(X \le x)$: the probability of a value at
    most $x$. \\
  $X \sim \dots$ & ``$X$ is distributed as'' & Names the distribution
    $X$ follows. \\
  $\E[X]$ & ``expected value of $X$'' & The mean: the probability-weighted
    average value. \\
  $\operatorname{Var}(X)$ & ``variance of $X$'' & The average squared
    distance from the mean: how spread out $X$ is. \\
  $\mu$, $\sigma$, $\sigma^2$ & ``mu, sigma, sigma squared'' & Mean,
    standard deviation, variance. \\
  $\mathcal{N}(\mu, \sigma^2)$ & ``normal with mean $\mu$ \dots'' & The
    Gaussian (bell curve) distribution. \\
  $\mathcal{N}(\vect{\mu}, \Sigma)$ & ``multivariate normal'' & A Gaussian
    over vectors, with mean vector $\vect{\mu}$ and covariance matrix
    $\Sigma$. \\
  $\Sigma$ & ``Sigma'' & The covariance matrix: variances on the diagonal,
    how coordinates move together elsewhere. \\
  $\text{Beta}(a, b)$ & ``Beta'' & A distribution over numbers between $0$
    and $1$, used as a prior for a probability. \\
  i.i.d. & ``independent and identically distributed'' & Data points drawn
    separately from the same distribution. \\
  $H(X)$, $H(\vect{p})$ & ``entropy of $X$'' & The average surprise
    $-\log p(x)$ of a discrete random variable: how uncertain it is. \\
  $h(X)$ & ``differential entropy'' & The entropy of a continuous random
    variable, computed from its density. \\
  $\E_{X \sim p}[\,\cdot\,]$ & ``expectation with $X$ drawn from $p$'' &
    An average in which $X$ follows the distribution $p$. \\
  $H(p, q)$ & ``cross-entropy of $q$ relative to $p$'' & The average
    surprise $-\log q(x)$ when $x$ is drawn from $p$: the cost of
    modeling $p$ with $q$. \\
  $D_{\mathrm{KL}}(p \,\|\, q)$ & ``KL divergence of $q$ from $p$'' &
    $H(p, q) - H(p)$: how much worse the model $q$ is than the truth
    $p$. Zero only when $q = p$; not symmetric. \\
  $I(X; Y)$ & ``mutual information'' & How much knowing $Y$ tells about
    $X$, and vice versa. \\
}

\section*{Statistics and Learning}

\symtable{
  $\mathcal{D}$ & ``the data'' & The observed data set. \\
  $N$ & ``N'' & The number of data points. \\
  $\theta$ & ``theta'' & An unknown parameter we want to learn. \\
  $\hat\theta$ & ``theta hat'' & An estimate of $\theta$ computed from
    data. \\
  $\bar{x}$ & ``$x$ bar'' & The sample mean, the average of the data. \\
  $L(\theta)$ & ``likelihood'' & How probable the observed data are if the
    parameter is $\theta$. \\
  $\ell(\theta)$ & ``log-likelihood'' & $\log L(\theta)$: the same
    information, as a sum instead of a product. \\
  $\hat\theta_{\text{MLE}}$, $\hat\theta_{\text{MAP}}$ & ``theta MLE, theta
    MAP'' & The maximum likelihood and maximum a posteriori estimates. \\
  $p(\theta)$, $p(\theta \mid \mathcal{D})$ & ``prior, posterior'' & What we
    believe about $\theta$ before and after seeing the data. \\
  $\vect{w}$, $b$ & ``weights, bias'' & The learned parameters of a model
    such as $\vect{w}^\top\vect{x} + b$. \\
  $y$, $\hat{y}$ & ``$y$, $y$ hat'' & The true label or target; the model's
    prediction. \\
  $\varepsilon$ & ``epsilon'' & Random noise, the part of $y$ the model
    cannot predict. \\
  $\lambda$ & ``lambda'' & The strength of a regularization penalty. \\
}

\section*{Greek Letters}

Greek letters are just names, like variable names in code. These are
the ones the book uses, with their usual job.

\symtable{
  $\alpha$ & alpha & An angle, such as the angle between two vectors. \\
  $\Delta$ & delta (capital) & A small change, as in $\Delta x$; $\Delta_n$
    is the simplex. \\
  $\varepsilon$ & epsilon & Random noise. \\
  $\eta$ & eta & Learning rate. \\
  $\theta$ & theta & A parameter; a probability of heads. \\
  $\lambda$, $\Lambda$ & lambda & Eigenvalue; rate of an exponential
    distribution; regularization strength. $\Lambda$ is a diagonal matrix
    of eigenvalues. \\
  $\mu$ & mu & Mean. \\
  $\pi$ & pi & The constant $3.14159\ldots$ \\
  $\sigma$, $\Sigma$ & sigma & Standard deviation; the sigmoid function;
    a singular value. $\Sigma$ is a covariance matrix, the middle factor
    of the SVD, or, written large as $\sum$, a sum. \\
  $\tau$ & tau & A standard deviation, usually of a prior. \\
  $\omega$, $\Omega$ & omega & An outcome; the sample space. \\
}

\begin{note}
  Some symbols have more than one job, as names do in code that reuses
  short variable names in different functions. The most common are
  $\sigma$ (standard deviation, sigmoid, singular value), $\Sigma$
  (covariance, SVD factor, sum), and $\lambda$ (eigenvalue, rate,
  regularization strength). The surrounding text always makes clear
  which one is meant.
\end{note}
\flushbottom
